\documentclass{amsart}

\usepackage[margin=1in]{geometry}
\usepackage[width=\linewidth]{caption}
\usepackage{tikz}

\definecolor{darkblue}{rgb}{0,0,0.4} 
 \usepackage[colorlinks=true, citecolor=darkblue, filecolor=darkblue, linkcolor=darkblue,urlcolor=darkblue]{hyperref} 
\usepackage[all]{hypcap}
\usepackage{xr-hyper}


\usepackage{ifthen,calc,amsmath,amssymb,mathtools,enumitem}
\usepackage{sseq}

\usepackage{xcolor,amssymb}
\usepackage{wasysym}
\usepackage{pdflscape}

\usepackage{pdfpages}

\definecolor{mygreen}{rgb}{0.09, 0.45, 0.27}

\newtheorem*{problem*}{Problem}

\def\Q {{\mathbb{Q}}}
\def\C{\mathbb{C}}
\def\R{\mathbb{R}}
\def\Z {\mathbb{Z}}
\def\Tor{\operatorname{Tor}}
\def\RP{\mathbb{RP}}
\def\KO{\mathit{KO}}
\def\ko{\mathit{ko}}
\def\Filt{\mathcal{F}}
\def\Gilt{\mathcal{G}}
\hbadness=100000

%SSmacros
\newcommand{\from}{\colon}
\newcommand{\Ext}{\mathrm{Ext}}
\newcommand{\Steenrod}{\mathcal{A}}
\newcommand{\Aone}{\Steenrod(1)}
\newcommand{\Azero}{\Steenrod(0)}
\newcommand{\FF}{\mathbb{F}}
\newcommand{\ZZ}{\mathbb{Z}}
\newcommand{\smas}{\wedge}
\newcommand{\Sq}{\mathrm{Sq}}
\newcommand{\defeq}{\coloneqq}


\begin{document}


\begin{problem*}
  For $n \geq 1$, consider the double cover $\pi\from S^n \to \RP^n$. Is image of the reduced $\KO$-theory map
  \[
  \pi^*\from \tilde{\KO}^{-1}(\RP^n) \to \tilde{\KO}^{-1}(S^n)
  \]
  $2$-divisible?
\end{problem*}


\begin{proof}[Solution]
  The image is indeed $2$-divisible.


  We will use $H^*$ to denote reduced cohomology, and $[\,]$ to denote
  homological grading shift. We write $\FF_2$ for $\ZZ/2$, and $\ZZ_2$
  for the $2$-adic integers. The tensor products are over $\Z$ unless
  otherwise noted. The Steenrod algebra over $\FF_2$ is $\Steenrod$;
  the subalgebra generated by $\Sq^1$ is
  $\Azero=\FF_2[\Sq^1]/(\Sq^1\Sq^1)$, and the subalgebra generated by
  $\Sq^1$ and $\Sq^2$ is
  \[
    \Aone = \FF_2\{\Sq^1, \Sq^2\} / (\Sq^1 \Sq^1=0, \ \Sq^1 \Sq^2 \Sq^1 = \Sq^2 \Sq^2),
  \]
  where the braces denote the free (non-commutative) algebra on those generators. 

  
  Let $D$ denote the Spanier-Whitehead dual; let $\RP^{-n}$ denote
  $D\RP^n$ and $S^{-n}$ denote $DS^n$. There are canonical
  isomorphisms $\KO^*(X)\cong \KO_{-*}(DX)$ such that
  \[
    \begin{tikzpicture}[xscale=3]
      \node (dr) at (0,0) {$\KO_{1}(\RP^{-n})$};
      \node (ds) at (1,0) {$\KO_{1}(S^{-n})$};
      \node (r) at (0,1) {$\KO^{-1}(\RP^n)$};
      \node (s) at (1,1) {$\KO^{-1}(S^n)$};
      
      \draw[->] (dr) -- (ds) node[pos=0.5, anchor=south] {\small $(D\pi)_*$};
      \draw[->] (r) -- (s) node[pos=0.5, anchor=south] {\small $\pi^*$};
      \draw[->] (r) -- (dr) node[pos=0.5,anchor=east] {\small $\cong$};
      \draw[->] (s) -- (ds) node[pos=0.5,anchor=west] {\small $\cong$};
    \end{tikzpicture}
  \]
  commutes.  Let $ko$ be the connective cover of the $\KO$
  spectrum. The spectra $X = S^{-n}$ and $X=\RP^{-n}$ are constructed
  from cells of non-positive dimension, so in both cases
  $\KO_1(X) \cong ko_1(X)$. Therefore, it is enough to prove for any
  $n\geq 1$, the image of
  \begin{equation}\label{eq:dual-main}
    (D\pi)_*\from ko_1(\RP^{-n}) \to ko_1(S^{-n})
  \end{equation}
  is $2$-divisible.


  
Let $X$ and $Y$ be CW spectra of finite type.
The Adams spectral sequence \cite{Adams, Brinkmann} (see also \cite{RognesNotes})
has the form
\begin{equation}
  E_2^{s, t}=\Ext^{s,t}_\Steenrod(H^*(Y;\FF_2),H^*(X;\FF_2))\Rightarrow   E_{\infty}^{t-s} = [X, Y]_{t-s} \otimes \ZZ_2,
\end{equation}
where $[X, Y]_k = [\Sigma^k X, Y]$ denotes the group of stable
homotopy classes of maps from $\Sigma^k X$ to $Y$. The differential
$d_r$ in the spectral sequence has bidegree $(r, r-1)$. Recall that
$ko_*(X)=[S,\ko \smas X]_*$. Therefore, we have an Adams spectral
sequence of the form
\[
\Ext^{s,t}_\Steenrod(H^*(\ko \smas X;\FF_2),\FF_2)\Rightarrow ko_{t-s}(X) \otimes \ZZ_2.
\]
 
A result of Stong \cite{Stong} implies that $H^*(\ko;\FF_2)\cong \Steenrod\otimes_{\Aone}\FF_2$ as
$\Steenrod$-modules; see \cite[Proposition 6.27]{RognesNotes} or \cite[Proposition 16.6]{Adams}. Therefore, 
\[
  H^*(\ko\smas X;\FF_2)\cong H^*(\ko;\FF_2)\otimes_{\FF_2}
  H^*(X;\FF_2)\cong \Steenrod\otimes_{\Aone} H^*(X;\FF_2),
\]
and hence,
\[
  \Ext^{s,t}_\Steenrod(H^*(\ko\smas X;\FF_2),\FF_2)\cong \Ext^{s,t}_{\Aone}(H^*(X;\FF_2),\FF_2).
\]
The Adams spectral sequence now reduces to:
\begin{equation}\label{eq:adams-sseq}
  \Ext^{s,t}_{\Aone}(H^*(X;\FF_2),\FF_2)\Rightarrow ko_{t-s}(X) \otimes \ZZ_2.
\end{equation}
When drawing this sequence, we will use the coordinate $i\defeq t-s$ on
the horizontal axis, and $s$ on the vertical axis; the differential $d_r$ changes the coordinates
$(i, s)$ by $(-1, r)$.



\begin{figure}
  \centering
\begin{sseq}[entrysize=5mm, gridstroke=.05pt]{0...12}{0...7}
\ssdropbull
\ssdroplabel{1}
\ssmove  0 1
\ssdropbull \ssstroke
\ssdroplabel{h_0}
\ssmove  0 1
\ssdropbull \ssstroke
\ssmove  0 1
\ssdropbull \ssstroke
\ssmove  0 1
\ssdropbull \ssstroke
\ssmove  0 1
\ssdropbull \ssstroke
\ssmove  0 1
\ssdropbull \ssstroke
\ssmove  0 1
\ssdropbull \ssstroke
\ssmove  0 1
\ssdropbull \ssstroke
\ssmove  0 1
\ssdropbull \ssstroke
\ssmove{0}{1}
\ssdropbull \ssstroke
\ssmove{0}{-10}
\ssline{1}{1}
\ssdropbull
\ssdroplabel{h_1}
\ssmove 1 1 
\ssdropbull \ssstroke
\ssmove 2 1
\ssdropbull
\ssdroplabel{v}
\ssmove  0 1
\ssdropbull \ssstroke
\ssmove  0 1
\ssdropbull \ssstroke
\ssmove  0 1
\ssdropbull \ssstroke
\ssmove  0 1
\ssdropbull \ssstroke
\ssmove  0 1
\ssdropbull \ssstroke
\ssmove{0}{1}
\ssdropbull \ssstroke
\ssmove{0}{-6}
\ssmove 4 1
\ssdropbull
\ssdroplabel{w_1}
\ssmove  0 1
\ssdropbull \ssstroke
%\ssdroplabel{h_0w_1}
\ssmove  0 1
\ssdropbull \ssstroke
\ssmove  0 1
\ssdropbull \ssstroke
\ssmove  0 1
\ssdropbull \ssstroke
\ssmove  0 1
\ssdropbull \ssstroke
\ssmove{0}{-5}
\ssline{1}{1}
\ssdropbull
%\ssdroplabel{h_1 \! w_1}
\ssmove 1 1 
\ssdropbull \ssstroke
\ssmove 2 1
\ssdropbull
%\ssdroplabel{vw_1}
\ssmove  0 1
\ssdropbull \ssstroke
\ssmove  0 1
\ssdropbull \ssstroke
\end{sseq}\qquad
\begin{sseq}[entrysize=5mm, gridstroke=.05pt]{0...12}{0...7}
  \ssdropbull
\ssmove  1 1
\ssdropbull \ssstroke
\ssmove  1 1
\ssdropbull \ssstroke
\ssmove  {0} {-1}
\ssdropbull \ssstroke
\ssmove  1 1
\ssdropbull \ssstroke
\ssmove  1 1
\ssdropbull \ssstroke
\ssmove 4 1
\ssdropbull
\ssmove  1 1
\ssdropbull \ssstroke
\ssmove  1 1
\ssdropbull \ssstroke
\ssmove  {0} {-1}
\ssdropbull \ssstroke
\ssmove  1 1
\ssdropbull \ssstroke
\ssmove  1 1
\ssdropbull \ssstroke
\end{sseq}
\caption{The Adams spectral sequence in $ko$-theory, for $S$ (left) and $\Sigma^{-1}\RP^2$ (right). Each is a module over $\Ext^{*,*}_{\Aone}(\FF_2,\FF_2)$; the vertical segments denote action of $h_0$ and the diagonal segments denote action of $h_1$. The patterns repeat infinitely on the right, after shifting gradings by $(8,4)$.}
\label{fig:spectralS}
\end{figure}


For $X=S$ the sphere spectrum, the spectral sequence is worked out
for example in \cite[Theorems 3.1.25 and 3.1.26]{Ravenel} or
\cite[Proposition 6.29]{RognesNotes}. Since
$H^*(S;\FF_2)=\FF_2$ in degree zero, the $E_2$ page is given by
\[
  E_2^{*, *} \cong \Ext^{*, *}_{\Aone}(\FF_2, \FF_2) \cong \FF_2[h_0, h_1, v, w_1] / (h_0h_1, h_1^3, h_1v, v^2=h_0^2w_1),
\]
with $h_0 \in E_2^{1,1}$, $h_1 \in E_2^{1,2}$, $v \in E_2^{3,7}$ and $w_1 \in E_2^{4, 12}$. The Adams spectral sequence in this case collapses at the $E_2$ page. After solving the extension problem, one arrives at (the $2$-adic completion of) the calculation of $ko$-theory:
\[
  ko_*(S) \otimes \ZZ_2 = \ZZ_2[\eta, \alpha, \beta]/(2\eta, \eta^3, \eta \alpha, \alpha^2 = 4\beta)
\]
where the classes $2, \eta, \alpha$ and $\beta$ came from
$h_0, h_1, v$ and $w_1$, respectively. Note that each infinite
vertical tower is the associated graded of a filtration on $\ZZ_2$. We
can read the groups $ko_*(S) \otimes \ZZ_2$ from each column:
$\ZZ_2, \ \ZZ/2, \ \ZZ/2, \ 0,\ \ZZ_2, \ 0, \ 0, \ 0,$ afterwards
repeated with period $8$.  More generally, if $X=S^{-n}$, then
$H^*(X;\FF_2)=\FF_2[n]$, and the spectral sequence in this case is the
same as above, but with the horizontal grading shifted by $-n$. Note
that for Equation~\eqref{eq:dual-main} we are interested in the
group $ko_1(S^{-n})= ko_{n+1}(S)$, which is only nonzero in degrees
\begin{equation}
\label{eq:ns}
 n \equiv 0, 1, 3, 7 \!\!\!\!\pmod{8}.
\end{equation}
When $n \equiv 0, 1 \!\!\pmod{8}$, the group $ko_1(S^{-n})$ is
$\ZZ/2$, and it remains $\ZZ/2$ after tensoring with $\ZZ_2$. Thus, we
want to show that the map $(D\pi)_*$ is zero, and it suffices to show
it is zero after tensoring with $\ZZ_2$. When
$n \equiv 3, 7 \!\!\pmod{8}$, the group $ko_1(S^{-n})$ is $\ZZ$. We
want to show that $(D\pi)_*$ has image in $2\ZZ \subset \ZZ$, and once
again it suffices to show that after tensoring with $\ZZ_2$, we land
in $2\ZZ_2 \subset \ZZ_2$.

Next, we consider the spectral sequence for $X=\Sigma^{-1}\RP^2$. Since
$H^*(\Sigma^{-1}\RP^2; \FF_2)\cong\Azero$, the $E_2$ term of the Adams spectral sequence is simply
$\Ext^{*, *}_{\Aone}(\Azero, \FF_2)$, which appears in \cite[Figure
24]{RognesNotes}. This spectral sequence also collapses at the $E_2$
page. These two spectral sequence are shown in
Figure~\ref{fig:spectralS}.



To prove our result, we need to understand
$ko_1(\RP^{-n}) \otimes \ZZ_2$. The action of $\Aone$ on
$H^*(\RP^{-n} ; \FF_2)$ is well-known, and is given by
\begin{equation}
\vcenter{\hbox{\begin{tikzpicture}[xscale=-1]
    \foreach\i in {1,...,8}{
      \node at (\i,0) {$\bullet$};
      \node[anchor=south] at (\i,0.1) {$-\i$};
    }
    \foreach \i in {1,3,...,7}{
      \draw (\i,0) -- (\i+1,0);
    }
    \foreach \i/\r in {2/1,3/-1,6/-1,7/1}{
      \draw (\i,0) to[out=\r*60,in=\r*120] (\i+2,0);
    }
    \node at (9.5,0) {$\cdots$};
  \end{tikzpicture}}}
\end{equation}
The tower continues periodically mod $4$, until it is terminated in degree $-n$.


To get a grasp on the $E_2$ page of the Adams spectral sequence \eqref{eq:adams-sseq} for $X=\RP^{-n}$, we will introduce two filtrations on $H^*(\RP^{-n}; \FF_2)$. We will study the induced filtrations on the $E_2$ page, and their associated graded modules.



The first filtration $\Filt$ on $H^*(\RP^{-n}; \FF_2)$ assigns to the generator $x_m \in H^m(\RP^{-n}; \FF_2)$ the filtration degree $\Filt(x_m) = k$ for  $m \in\{  -2k-2, -2k-1\}$; $\Sq^1$ preserves the filtration degree, and $\Sq^2$ decreases it, and the associated graded is given by
\begin{equation}
\vcenter{\hbox{\begin{tikzpicture}[xscale=-1]
    \foreach\i in {1,...,8}{
      \node at (\i,0) {$\bullet$};
    }
    \foreach \i in {0,1,2,3}{
      \draw (2*\i+1,0) -- (2*\i+2,0);
      \node[anchor=south] at (2*\i+1.5,0) {$F_\i$};
    }
    \node at (9.5,0) {$\cdots$};
  \end{tikzpicture}}}
\end{equation}
where $F_k$ denotes the associated graded of $\Filt$ in filtration
degree $k$. It is usually isomorphic to $\Azero[-2k]$, except that when $n$
is odd the last piece $F_{(n-1)/2}$ is just a copy of $\FF_2[-n]$. Therefore, the associated graded
of the $E_2$ page of the Adams spectral sequence has $n/2$ shifted
copies of the one for $\Sigma^{-1}\RP^2$ starting in degrees
$-n, -n+2, \dots, -2$ for $n$ even, and has $(n-1)/2$ shifted copies
of the one for $\Sigma^{-1}\RP^2$ starting in degrees
$-n+1, -n+3, \dots, -2$ and one copy of the one for $S$ starting in
degree $-n$, for $n$ odd.  See Figure~\ref{fig:Filt}.


\begin{figure}
  \centering
\begin{sseq}[entrysize=5mm, gridstroke=.05pt, xlabels={-n; ; ;;;;; -; ;;; ; -n+12}]{0...12}{0...7}
\ssdropbull
\ssmove  1 1
\ssdropbull \ssstroke
\ssmove  1 1
\ssdropbull \ssstroke
\ssmove  {0} {-1}
\ssdropbull \ssstroke
\ssmove  1 1
\ssdropbull \ssstroke
\ssmove  1 1
\ssdropbull \ssstroke
\ssmove 4 1
\ssdropbull
\ssmove  1 1
\ssdropbull \ssstroke
\ssmove  1 1
\ssdropbull \ssstroke
\ssmove  {0} {-1}
\ssdropbull \ssstroke
\ssmove  1 1
\ssdropbull \ssstroke
\ssmove  1 1
\ssdropbull \ssstroke
\ssmove {-10}{-7}
\ssdropbull
\ssmove  1 1
\ssdropbull \ssstroke
\ssmove  1 1
\ssdropbull \ssstroke
\ssmove  {0} {-1}
\ssdropbull \ssstroke
\ssmove  1 1
\ssdropbull \ssstroke
\ssmove  1 1
\ssdropbull \ssstroke
\ssmove 4 1
\ssdropbull
\ssmove  1 1
\ssdropbull \ssstroke
\ssmove  1 1
\ssdropbull \ssstroke
\ssmove  {0} {-1}
\ssdropbull \ssstroke
\ssmove  1 1
\ssdropbull \ssstroke
\ssmove  1 1
\ssdropbull \ssstroke
\ssmove {-10}{-7}
\ssdropbull
\ssmove  1 1
\ssdropbull \ssstroke
\ssmove  1 1
\ssdropbull \ssstroke
\ssmove  {0} {-1}
\ssdropbull \ssstroke
\ssmove  1 1
\ssdropbull \ssstroke
\ssmove  1 1
\ssdropbull \ssstroke
\ssmove 4 1
\ssdropbull
\ssmove  1 1
\ssdropbull \ssstroke
\ssmove  1 1
\ssdropbull \ssstroke
\ssmove  {0} {-1}
\ssdropbull \ssstroke
\ssmove  1 1
\ssdropbull \ssstroke
\ssmove  1 1
\ssdropbull \ssstroke
\ssmove {-10}{-7}
\ssdropbull
\ssmove  1 1
\ssdropbull \ssstroke
\ssmove  1 1
\ssdropbull \ssstroke
\ssmove  {0} {-1}
\ssdropbull \ssstroke
\ssmove  1 1
\ssdropbull \ssstroke
\ssmove  1 1
\ssdropbull \ssstroke
\ssmove 4 1
\ssdropbull
\ssmove  1 1
\ssdropbull \ssstroke
\ssmove  1 1
\ssdropbull \ssstroke
\ssmove  {0} {-1}
\ssdropbull \ssstroke
\ssmove  1 1
\ssdropbull \ssstroke
\ssmove  1 1
\ssdropbull \ssstroke
\ssmove {-10}{-7}
\ssdropbull
\ssmove  1 1
\ssdropbull \ssstroke
\ssmove  1 1
\ssdropbull \ssstroke
\ssmove  {0} {-1}
\ssdropbull \ssstroke
\ssmove  1 1
\ssdropbull \ssstroke
\ssmove  1 1
\ssdropbull \ssstroke
\ssmove 4 1
\ssdropbull
\ssmove  1 1
\ssdropbull \ssstroke
\ssmove  1 1
\ssdropbull \ssstroke
\ssmove  {0} {-1}
\ssdropbull \ssstroke
\ssmove  1 1
\ssdropbull \ssstroke
\ssmove  1 1
\ssdropbull \ssstroke
\ssmove {-10}{-7}
\ssdropbull
\ssmove  1 1
\ssdropbull \ssstroke
\ssmove  1 1
\ssdropbull \ssstroke
\ssmove  {0} {-1}
\ssdropbull \ssstroke
\ssmove  1 1
\ssdropbull \ssstroke
\ssmove  1 1
\ssdropbull \ssstroke
\ssmove 4 1
\ssdropbull
\ssmove  1 1
\ssdropbull \ssstroke
\ssmove  1 1
\ssdropbull \ssstroke
\ssmove  {0} {-1}
\ssdropbull \ssstroke
\ssmove  1 1
\ssdropbull \ssstroke
\ssmove  1 1
\ssdropbull \ssstroke
\ssmove {-10}{-7}
\ssdropbull
\ssmove  1 1
\ssdropbull \ssstroke
\ssmoveto 2 2
\ssline[dashed]{12}{6}
\ssdropbull
\ssmoveto 2 2
\ssline[dashed]{-4}{-2}
\ssdropbull
\end{sseq}\qquad
\begin{sseq}[entrysize=5mm, gridstroke=.05pt, xlabels={-n; ; ;; ; ; ; ; ;; ; ; -n+12}]{0...12}{0...7}
\ssdropbull
\ssmove  0 1
\ssdropbull \ssstroke
\ssmove  0 1
\ssdropbull \ssstroke
\ssmove  0 1
\ssdropbull \ssstroke
\ssmove  0 1
\ssdropbull \ssstroke
\ssmove  0 1
\ssdropbull \ssstroke
\ssmove  0 1
\ssdropbull \ssstroke
\ssmove  0 1
\ssdropbull \ssstroke
\ssmove  0 1
\ssdropbull \ssstroke
\ssmove  0 1
\ssdropbull \ssstroke
\ssmove{0}{1}
\ssdropbull \ssstroke
\ssmove{0}{-10}
\ssline{1}{1}
\ssdropbull
\ssmove 1 1 
\ssdropbull \ssstroke
\ssmove 2 1
\ssdropbull
\ssmove  0 1
\ssdropbull \ssstroke
\ssmove  0 1
\ssdropbull \ssstroke
\ssmove  0 1
\ssdropbull \ssstroke
\ssmove  0 1
\ssdropbull \ssstroke
\ssmove  0 1
\ssdropbull \ssstroke
\ssmove {0}{-5}
\ssmove 4 1
\ssdropbull
\ssmove  0 1
\ssdropbull \ssstroke
%\ssdroplabel{h_0w_1}
\ssmove  0 1
\ssdropbull \ssstroke
\ssmove  0 1
\ssdropbull \ssstroke
\ssmove  0 1
\ssdropbull \ssstroke
\ssmove  0 1
\ssdropbull \ssstroke
\ssmove{0}{-5}
\ssline{1}{1}
\ssdropbull
%\ssdroplabel{h_1 \! w_1}
\ssmove 1 1 
\ssdropbull \ssstroke
\ssmove 2 1
\ssdropbull
\ssmove  0 1
\ssdropbull \ssstroke
\ssmove  0 1
\ssdropbull \ssstroke

\ssmoveto 1 0
\ssdropbull
\ssmove  1 1
\ssdropbull \ssstroke
\ssmove  1 1
\ssdropbull \ssstroke
\ssmove  {0} {-1}
\ssdropbull \ssstroke
\ssmove  1 1
\ssdropbull \ssstroke
\ssmove  1 1
\ssdropbull \ssstroke
\ssmove 4 1
\ssdropbull
\ssmove  1 1
\ssdropbull \ssstroke
\ssmove  1 1
\ssdropbull \ssstroke
\ssmove  {0} {-1}
\ssdropbull \ssstroke
\ssmove  1 1
\ssdropbull \ssstroke
\ssmove  1 1
\ssdropbull \ssstroke
\ssmove {-10}{-7}
\ssdropbull
\ssmove  1 1
\ssdropbull \ssstroke
\ssmove  1 1
\ssdropbull \ssstroke
\ssmove  {0} {-1}
\ssdropbull \ssstroke
\ssmove  1 1
\ssdropbull \ssstroke
\ssmove  1 1
\ssdropbull \ssstroke
\ssmove 4 1
\ssdropbull
\ssmove  1 1
\ssdropbull \ssstroke
\ssmove  1 1
\ssdropbull \ssstroke
\ssmove  {0} {-1}
\ssdropbull \ssstroke
\ssmove  1 1
\ssdropbull \ssstroke
\ssmove  1 1
\ssdropbull \ssstroke
\ssmove {-10}{-7}
\ssdropbull
\ssmove  1 1
\ssdropbull \ssstroke
\ssmove  1 1
\ssdropbull \ssstroke
\ssmove  {0} {-1}
\ssdropbull \ssstroke
\ssmove  1 1
\ssdropbull \ssstroke
\ssmove  1 1
\ssdropbull \ssstroke
\ssmove 4 1
\ssdropbull
\ssmove  1 1
\ssdropbull \ssstroke
\ssmove  1 1
\ssdropbull \ssstroke
\ssmove  {0} {-1}
\ssdropbull \ssstroke
\ssmove  1 1
\ssdropbull \ssstroke
\ssmove  1 1
\ssdropbull \ssstroke
\ssmove {-10}{-7}
\ssdropbull
\ssmove  1 1
\ssdropbull \ssstroke
\ssmove  1 1
\ssdropbull \ssstroke
\ssmove  {0} {-1}
\ssdropbull \ssstroke
\ssmove  1 1
\ssdropbull \ssstroke
\ssmove  1 1
\ssdropbull \ssstroke
\ssmove 4 1
\ssdropbull
\ssmove  1 1
\ssdropbull \ssstroke
\ssmove  1 1
\ssdropbull \ssstroke
\ssmove  {0} {-1}
\ssdropbull \ssstroke
\ssmove  1 1
\ssdropbull \ssstroke
\ssmove  1 1
\ssdropbull \ssstroke
\ssmove {-10}{-7}
\ssdropbull
\ssmove  1 1
\ssdropbull \ssstroke
\ssmove  1 1
\ssdropbull \ssstroke
\ssmove  {0} {-1}
\ssdropbull \ssstroke
\ssmove  1 1
\ssdropbull \ssstroke
\ssmove  1 1
\ssdropbull \ssstroke
\ssmove 4 1
\ssdropbull
\ssmove  1 1
\ssdropbull \ssstroke
\ssmove  1 1
\ssdropbull \ssstroke
\ssmove  {0} {-1}
\ssdropbull \ssstroke
\ssmove  1 1
\ssdropbull \ssstroke
\ssmove  1 1
\ssdropbull \ssstroke
\ssmove {-10}{-7}
\ssdropbull
\ssmove  1 1
\ssdropbull \ssstroke
\ssmove  1 1
\ssdropbull \ssstroke
\ssmove  {0} {-1}
\ssdropbull \ssstroke
\ssmove  1 1
\ssdropbull \ssstroke
\ssmove  1 1
\ssdropbull \ssstroke
\ssmove 4 1
\ssdropbull
\ssmove  1 1
\ssdropbull \ssstroke
\ssmove  1 1
\ssdropbull \ssstroke
\ssmove  {0} {-1}
\ssdropbull \ssstroke
\ssmove  1 1
\ssdropbull \ssstroke
\ssmove  1 1
\ssdropbull \ssstroke
\ssmove {-10}{-7}
\ssdropbull
\ssmove  1 1
\ssdropbull \ssstroke
\ssmoveto 1 1
\ssline[dashed]{14}{7}
\ssdropbull
\ssmoveto 1 1
\ssline[dashed]{-4}{-2}
\ssdropbull
\end{sseq}
\caption{The $\Filt$-associated graded of the $E_2$ page for $n$ even (left) and $n$ odd (right).}
\label{fig:Filt}
\end{figure}


The actual page
$E_2 = \Ext_{\Aone}^{*, *}(H^*(\RP^{-n}; \FF_2), \FF_2)$ is obtained
from the associated graded by canceling various generators related by
differentials coming from $\Filt$. We do not yet know what the
differentials are, but in any case, for $n$ even, the result will be
supported in the half-plane
\begin{equation}\label{eq:even-support}
  s\leq\frac{i}{2}+ \frac{n}{2} + 1,
\end{equation}
and
for $n$ odd, the parts of the $E_2$ page supported in columns with $i$
even lie in the half-plane
\begin{equation}\label{eq:odd-support}
  s \leq \frac{i}{2} + \frac{n+1}{2};
\end{equation}
see
Figure~\ref{fig:Filt}.


The second filtration $\Gilt$ assigns to $x_m \in H^m(\RP^{-n}; \FF_2)$ the filtration degree $\Gilt(x_m) = k$  for  $m \in\{  4k-4, 4k-2, 4k-1, 4k+1\};$
$\Sq^2$ preserves the filtration degree and $\Sq^1$ does not increase it,
and the associated graded is given by
\begin{equation}
\vcenter{\hbox{\begin{tikzpicture}[xscale=-1]
    \foreach\i in {1,...,8}{
      \node at (\i,0) {$\bullet$};
    }
    \foreach \i in {0,2}{
      \draw (2*\i+1,0) -- (2*\i+2,0);
    }
    \foreach \i/\r in {2/1,3/-1,6/-1,7/1}{
      \draw (\i,0) to[out=\r*60,in=\r*120] (\i+2,0);
    }
    \node at (9.5,0) {$\cdots$};
    \node[anchor=south] at (1.5,0) {$G_0$};
    \node[anchor=south] at (5.5,0) {$G_{-1}$};
    
  \end{tikzpicture}}}
\end{equation}
where $G_k$ denotes the associated graded of $\Gilt$ in filtration
degree $k$. Let $P\defeq G_{-1}[8]$ be the $4$-dimensional module
shown above, shifted so that it is supported in gradings $0,2,3,5$. In
general, most of the modules $G_k$ are isomorphic to (shifted copies
of) $P$; the exceptions are $G_0$ and those near the bottom of the
tower, which are truncations of $P$, and their shape depends on the
congruence type of $n$ mod $4$. Let $R_1$ denote the submodule of $P$
generated by elements in gradings $3$ and $5$, but shifted to start in
grading $0$. Let $R_2$ denote the submodule of $P$ generated by
elements in grading $2,3,5$, again shifted to start in grading
$0$. Then, for $n\equiv 0,1,2,3\!\!\pmod 4$, the lowest associated
graded pieces are $\FF_2[1-n]$, $R_1[-n]$, $R_2[-n]$, and
$\FF_2[-n]\oplus R_2[1-n]$, respectively. (For $n\equiv 3\!\!\pmod 4$,
$\FF_2[-n]$ and $R_2[1-n]$ are in different filtration gradings.)


From now on let us assume that $n \geq 4$, so that $G_0$ if of the
form described above. (The cases $n \leq 3$ will be dealt with
separately.) We would like to understand the $E_2$ pages of the
associated graded pieces, that is, $\Ext_{\Aone}^{s, t}(M, \FF_2)$,
for $M = \FF_2, P, G_0, R_1, R_2.$ The case $M=\FF_2$ is that of the
sphere spectrum, already discussed earlier.

For $M=P$, observe that $P \cong \Aone \otimes_{\Azero} \FF_2,$
and therefore
\[
  \Ext_{\Aone}^{s, t}(P, \FF_2) \; \cong  \; \Ext_{\Aone}^{s, t}( \Aone \otimes_{\Azero} \FF_2, \FF_2)   \;\cong  \; \Ext_{\Azero}^{s, t}(\FF_2, \FF_2) \;  \cong  \; \FF_2[x].
\]
We conclude that the associated graded piece $G_k=P[4k-4]$ yields a vertical tower $\FF_2[x]$ in the associated graded of the $E_2$ page of \eqref{eq:adams-sseq}; the tower has one generator in each bidegree $(i, s)=(4k-4, s)$ for $s \geq 0$.

Next, let us consider $M=R_1$. Let $E(1)$ be the subalgebra of $\Aone$ generated by $\Sq^1$ and the commutator $Q_1\defeq[\Sq^1, \Sq^2]$. Note that $E(1)$ is the exterior algebra on $\Sq^1$ and $Q_1$. We have
$ R_1 \cong H^*(ku) \cong  \Aone \otimes_{E(1)} \FF_2.$
Therefore
\[
  \Ext_{\Aone}^{s, t}(Q, \FF_2) \; \cong  \; \Ext_{\Aone}^{s, t}( \Aone \otimes_{E(1)} \FF_2, \FF_2)   \;\cong  \; \Ext_{E(1)}^{s, t}(\FF_2, \FF_2)  \;  \cong  \; \FF_2[h_0, v_1],
\]
where the last isomorphism comes from the K\"unneth formula and Koszul duality. This is in fact a standard calculation for the $E_2$ page of the Adams spectral sequence for $ku$; see for example \cite[Proposition 6.28]{RognesNotes}. The element $h_0$ lives in bidegree $(i, s)=(0,1)$ and $v_1$ in bidegree $(i, s)=(2,1)$.

\begin{figure}
  \centering
\begin{sseq}[entrysize=5mm, gridstroke=.05pt]{0...8}{0...5}
\ssdropbull
\ssdroplabel{1}
\ssmove  0 1
\ssdropbull \ssstroke
\ssdroplabel{h_0}
\ssmove  0 1
\ssdropbull \ssstroke
\ssmove  0 1
\ssdropbull \ssstroke
\ssmove  0 1
\ssdropbull \ssstroke
\ssmove  0 1
\ssdropbull \ssstroke
\ssmove  0 1
\ssdropbull \ssstroke
\ssmove  {2} {-5}
\ssdropbull
\ssdroplabel{v_1}
\ssmove  0 1
\ssdropbull \ssstroke
\ssmove  0 1
\ssdropbull \ssstroke
\ssmove  0 1
\ssdropbull \ssstroke
\ssmove  0 1
\ssdropbull \ssstroke
\ssmove  0 1
\ssdropbull \ssstroke
\ssmove  {2} {-4}
\ssdropbull
\ssmove  0 1
\ssdropbull \ssstroke
\ssmove  0 1
\ssdropbull \ssstroke
\ssmove  0 1
\ssdropbull \ssstroke
\ssmove  0 1
\ssdropbull \ssstroke
\ssmove{2}{-3}
\ssdropbull
\ssmove  0 1
\ssdropbull \ssstroke
\ssmove  0 1
\ssdropbull \ssstroke
\ssmove  0 1
\ssdropbull \ssstroke
\ssmove{2}{-2}
\ssdropbull
\ssmove  0 1
\ssdropbull \ssstroke
\ssmove  0 1
\ssdropbull \ssstroke
\end{sseq}\qquad
\begin{sseq}[entrysize=5mm, gridstroke=.05pt]{-4...4}{0...5}
\ssmoveto {-4}{0}
\ssdropbull
\ssmove  0 1
\ssdropbull \ssstroke
\ssmove  0 1
\ssdropbull \ssstroke
\ssmove  0 1
\ssdropbull \ssstroke
\ssmove  0 1
\ssdropbull \ssstroke
\ssmove  0 1
\ssdropbull \ssstroke
\ssmove  0 1
\ssdropbull \ssstroke
\ssmove  0 1
\ssdropbull \ssstroke
\ssmove  0 1
\ssdropbull \ssstroke
\ssmove  0 1
\ssdropbull \ssstroke
\ssmoveto{0}{1}
\ssdropbull
\ssmove  0 1
\ssdropbull \ssstroke
\ssmove  0 1
\ssdropbull \ssstroke
\ssmove  0 1
\ssdropbull \ssstroke
\ssmove  0 1
\ssdropbull \ssstroke
\ssmove  0 1
\ssdropbull \ssstroke
\ssmove  0 1
\ssdropbull \ssstroke
\ssmove  0 1
\ssdropbull \ssstroke
\ssmove  0 1
\ssdropbull \ssstroke
\ssmoveto {0}{1}
\ssline  1 1
\ssdropbull \ssstroke
\ssmove  1 1
\ssdropbull \ssstroke
\ssmoveto{4}{4}
\ssdropbull
\ssmove  0 1
\ssdropbull \ssstroke
\ssmove  0 1
\ssdropbull \ssstroke
\ssmove  0 1
\ssdropbull \ssstroke
\ssmove  0 1
\ssdropbull \ssstroke
\end{sseq}
\caption{The $\Ext$ groups for $R_1$ (left) and $G_0$ (right), both repeating with periodicity $(8,4)$.}
\label{fig:R1}
\end{figure}

For $M = G_0$, starting in degree $-4$, the chart is also well-known, appearing for example in \cite[Figure 12]{RognesNotes}, for the calculation of $H^*(b spin)$. (One can also calculate this directly using the techniques that we will employ for $R_2$ below.) Figure~\ref{fig:R1} shows the charts for $R_1$ and $G_0$.

Finally for $M=R_2=\vcenter{\hbox{\begin{tikzpicture}
      \foreach \i in {0,1,3}{\node at (\i,0) {$\bullet$};}
      \draw (0,0)--(1,0);
      \draw (1,0) to[out=60,in=120] (3,0);
    \end{tikzpicture}}}$, we may filter it in two ways with the following associated graded pieces:
\[
  \begin{tikzpicture}
    \begin{scope}[xshift=-5cm]
      \foreach \i in {0,1,3}{\node at (\i,0) {$\bullet$};}
      \draw (1,0) to[out=60,in=120] (3,0);
    \end{scope}
      \node at (-1,0) {\text{and}};
      \foreach \i in {0,1,3}{\node at (\i,0) {$\bullet$};}
      \draw (0,0)--(1,0);
    \end{tikzpicture}
\]
The first one is $\FF_2\oplus R_1[1]$ and the second one is
$\Azero\oplus\FF_2[3]$. Each gives a spectral sequence converging to
$\Ext_{\Aone}^{*,*}(R_2,\FF_2)$, and the differentials in one are
determined by the shape of the other. We show the associated graded of
the $E_2$ pages for the above two filtrations in Figure~\ref{fig:R2};
the differentials are shown in the figure and they are forced by the
two shapes (and the fact that they have to commute with the $h_0$ and
$h_1$ actions). The combined $\Ext$ group is also shown. The only
unknown is whether the dotted line in the figure (representing a
possible action of $h_1$) connects the generators in grading $(7,3)$
and $(8,4)$; we will not need to know this for our purposes.


\begin{figure}
  \centering
\begin{sseq}[entrysize=5mm, gridstroke=.05pt]{0...8}{0...6}
\ssdropbull
\ssmove  1 1
\ssdropbull \ssstroke[color=red]
\ssmove 1 1
\ssdropbull \ssstroke[color=red]
\ssmoveto{0}{0}
\ssline{0}{1}[color=red]
\ssdropbull \ssstroke[color=red]
\ssmove  0 1
\ssdropbull \ssstroke[color=red]
\ssmove  0 1
\ssdropbull \ssstroke[color=red]
\ssmove  0 1
\ssdropbull \ssstroke[color=red]
\ssmove  0 1
\ssdropbull \ssstroke[color=red]
\ssmove  0 1
\ssdropbull \ssstroke[color=red]
\ssmove  0 1
\ssdropbull \ssstroke[color=red]

\ssmoveto{4}{3}
\ssdropbull
\ssmove 0 1
\ssdropbull \ssstroke[color=red]
\ssmove 0 1
\ssdropbull \ssstroke[color=red]
\ssmove 0 1
\ssdropbull \ssstroke[color=red]
\ssmove 0 1
\ssdropbull \ssstroke[color=red]
\ssmove 0 1
\ssdropbull \ssstroke[color=red]
\ssmove 0 1
\ssdropbull \ssstroke[color=red]

\ssmoveto{8}{4}
\ssdropbull
\ssmove 1 1
\ssdropbull \ssstroke[color=red]
\ssmoveto{8}{4}
\ssline{0}{1}[color=red]
\ssdropbull \ssstroke[color=red]
\ssmove 0 1
\ssdropbull \ssstroke[color=red]
\ssmove 0 1
\ssdropbull \ssstroke[color=red]

\ssmoveto{1}{0}
\ssdropbull
\ssmove 0 1
\ssdropbull \ssname{blue} \ssstroke[color=blue]
\ssmove 0 1
\ssdropbull \ssstroke[color=blue]
\ssmove 0 1
\ssdropbull \ssstroke[color=blue]
\ssmove 0 1
\ssdropbull \ssstroke[color=blue]
\ssmove 0 1
\ssdropbull \ssstroke[color=blue]
\ssmove 0 1
\ssdropbull \ssstroke[color=blue]
\ssmove 0 1
\ssdropbull \ssstroke[color=blue]
\ssmove 0 1
\ssdropbull \ssstroke[color=blue]

\ssmoveto{3}{1}
\ssdropbull
\ssmove 0 1
\ssdropbull \ssstroke[color=blue]
\ssmove 0 1
\ssdropbull \ssstroke[color=blue]
\ssmove 0 1
\ssdropbull \ssstroke[color=blue]
\ssmove 0 1
\ssdropbull \ssstroke[color=blue]
\ssmove 0 1
\ssdropbull \ssstroke[color=blue]
\ssmove 0 1
\ssdropbull \ssstroke[color=blue]

\ssmoveto{5}{2}
\ssdropbull
\ssmove 0 1
\ssdropbull \ssstroke[color=blue]
\ssmove 0 1
\ssdropbull \ssstroke[color=blue]
\ssmove 0 1
\ssdropbull \ssstroke[color=blue]
\ssmove 0 1
\ssdropbull \ssstroke[color=blue]
\ssmove 0 1
\ssdropbull \ssstroke[color=blue]
\ssmove 0 1
\ssdropbull \ssstroke[color=blue]

\ssmoveto{7}{3}
\ssdropbull
\ssmove 0 1
\ssdropbull \ssstroke[color=blue]
\ssmove 0 1
\ssdropbull \ssstroke[color=blue]
\ssmove 0 1
\ssdropbull \ssstroke[color=blue]
\ssmove 0 1
\ssdropbull \ssstroke[color=blue]
\ssmove 0 1
\ssdropbull \ssstroke[color=blue]
\ssmove 0 1
\ssdropbull \ssstroke[color=blue]

\ssmoveto{3}{1}
\ssline[dashed]{-1}{1}\ssarrowhead

\ssmoveto{1}{0}
\ssline[dashed]{-1}{1}\ssarrowhead
\ssmoveto{1}{2}
\ssline[dashed]{-1}{1}\ssarrowhead
\ssmoveto{1}{3}
\ssline[dashed]{-1}{1}\ssarrowhead
\ssmoveto{1}{4}
\ssline[dashed]{-1}{1}\ssarrowhead
\ssmoveto{1}{5}
\ssline[dashed]{-1}{1}\ssarrowhead
\ssmoveto{1}{6}
\ssline[dashed]{-1}{1}\ssarrowhead

\ssgoto{blue}
\ssline[dashed, curve=-.2]{-1}{1}\ssarrowhead

\ssmoveto{5}{2}
\ssline[dashed]{-1}{1}\ssarrowhead
\ssmoveto{5}{3}
\ssline[dashed]{-1}{1}\ssarrowhead
\ssmoveto{5}{4}
\ssline[dashed]{-1}{1}\ssarrowhead
\ssmoveto{5}{5}
\ssline[dashed]{-1}{1}\ssarrowhead
\ssmoveto{5}{6}
\ssline[dashed]{-1}{1}\ssarrowhead

\ssmoveto{9}{4}
\ssdropbull \ssline[dashed]{-1}{1}\ssarrowhead
\ssmoveto{9}{5}
\ssdropbull \ssline[dashed]{-1}{1}\ssarrowhead

\end{sseq}\
\begin{sseq}[entrysize=5mm, gridstroke=.05pt,ylabels=none]{0...8}{0...6}

  \ssdropbull
  \ssmove 1 1
  \ssdropbull \ssstroke[color=red]
  \ssmove 1 1
  \ssdropbull \ssstroke[color=red]
  \ssmove {0}{-1}
  \ssdropbull \ssstroke[color=red]
  \ssmove 1 1
  \ssdropbull \ssname{red} \ssstroke[color=red]
  \ssmove 1 1
  \ssdropbull \ssstroke[color=red]

  \ssmoveto{8}{4}
  \ssdropbull
  \ssmove 1 1
  \ssdropbull \ssstroke[color=red]

  \ssmoveto{3}{0}
  \ssdropbull 
  \ssmove 0 1
  \ssdropbull \ssstroke[color=blue]
  \ssmove 0 1
  \ssdropbull \ssstroke[color=blue]
  \ssmove 0 1
  \ssdropbull \ssstroke[color=blue]
  \ssmove 0 1
  \ssdropbull \ssstroke[color=blue]
  \ssmove 0 1
  \ssdropbull \ssstroke[color=blue]
  \ssmove 0 1
  \ssdropbull \ssstroke[color=blue]
  \ssmove 0 1
  \ssdropbull \ssstroke[color=blue]
  \ssmove 0 1
  \ssdropbull \ssstroke[color=blue]
  \ssmoveto{3}{0}
  \ssline{1}{1}[color=blue]
  \ssdropbull \ssname{blue}
  \ssmove 1 1
  \ssdropbull \ssstroke[color=blue]

  \ssmoveto{7}{3}
  \ssdropbull
  \ssmove 0 1
  \ssdropbull \ssstroke[color=blue]
  \ssmove 0 1
  \ssdropbull \ssstroke[color=blue]
  \ssmove 0 1
  \ssdropbull \ssstroke[color=blue]
  \ssmove 0 1
  \ssdropbull \ssstroke[color=blue]
  \ssmove 0 1
  \ssdropbull \ssstroke[color=blue]
  \ssmove 0 1
  \ssdropbull \ssstroke[color=blue]

\ssmoveto{3}{0}
\ssline[dashed] {-1}{1}\ssarrowhead
\ssmoveto{5}{2}
\ssline[dashed] {-1}{1}\ssarrowhead
\ssmoveto{3}{1}
\ssline[dashed] {-1}{1}\ssarrowhead
\ssgoto{blue} \ssgoto{red}
\ssstroke[dashed,curve=-0.2] \ssarrowhead

  
\end{sseq}\
\begin{sseq}[entrysize=5mm, gridstroke=.05pt,ylabels=none]{0...8}{0...6}
\ssdropbull
\ssmove  1 1
\ssdropbull \ssstroke
\ssmoveto{3}{2}
\ssdropbull
\ssmove  0 1
\ssdropbull \ssstroke
\ssmove  0 1
\ssdropbull \ssstroke
\ssmove  0 1
\ssdropbull \ssstroke
\ssmove  0 1
\ssdropbull \ssstroke
\ssmove  0 1
\ssdropbull \ssstroke
\ssmove  0 1
\ssdropbull \ssstroke

\ssmoveto{7}{3}
\ssdropbull
\ssmove  0 1
\ssdropbull \ssstroke
\ssmove  0 1
\ssdropbull \ssstroke
\ssmove  0 1
\ssdropbull \ssstroke
\ssmove  0 1
\ssdropbull \ssstroke
\ssmoveto 7 3
\ssline[dotted]{1}{1}
\ssdropbull
\ssmove  1 1
\ssdropbull \ssstroke
\end{sseq}
\caption{The associated graded for the two filtrations, along with the
  differentials (left and center), and the $\Ext$ groups for $R_2$
  (right); for the associated graded, the parts coming from the two
  pieces are drawn in different colors. The picture repeats with
  periodicity $(8, 4)$.}
\label{fig:R2}
\end{figure}

We can now study $\ko_1(\RP^{-n})$. Let $n = 8j+k$ with
$0 \leq k \leq 7$. We consider the filtration $\Gilt$ on the $E_2$
page. Following the above discussion, we superimpose the following
charts:
\begin{itemize}[leftmargin=*]
\item the chart for $G_0$ (Figure~\ref{fig:R1} right);
\item $2j$ vertical lines, coming from the copies of $P$, in degrees $i \in \{-8j, -8j+4, \dots, -12, -8\}$;
\item depending on the congruence class of $n$ mod $4$, the chart for
  $\FF_2$ (Figure~\ref{fig:spectralS} left) shifted left by $n-1$, the
  chart for $R_1$ (Figure~\ref{fig:R1} left) shifted left by $n$, the
  chart for $R_2$ (Figure~\ref{fig:R2}), or the chart for $\FF_2$
  shifted left by $n$ and the chart for $R_2$ shifted left by $n-1$.
\end{itemize}

Since we are interested in $ko_1(\RP^{-n})$, we focus on the columns
in the chart near $i=1$. These are drawn in Figure~\ref{fig:n07}, for
each congruence class of $n$ mod $8$.  The red patterns come from
$G_0$, the green patterns from $\FF_2$ and the blue patterns from
either $R_1$ or $R_2$. For $n=8j+k$ with $k \in \{0,1,2,3\}$, we drew
the case $j=1$, and for $k \in \{4,5,6,7\}$, the case $j=0$. In
general, the blue and green patterns start in degrees shifted upwards
by multiples of $4$; for example, for $n=8j$, the green pattern starts
in bi-degree $(1, 4j)$.

 

\begin{figure}
   \centering
\begin{sseq}[entrysize=5mm, gridstroke=.05pt,xlabelstep=1]{0...2}{0...8}
\ssmove 0 1
\ssdropbull
\ssmove  0 1
\ssdropbull  \ssstroke[color=red]
\ssmove  0 1
\ssdropbull \ssstroke[color=red]
\ssmove  0 1
\ssdropbull \ssstroke[color=red]
\ssmove  0 1
\ssdropbull \ssstroke[color=red]
\ssmove  0 1
\ssdropbull \ssstroke[color=red]
\ssmove  0 1
\ssdropbull \ssstroke[color=red]
\ssmove  0 1
\ssdropbull \ssstroke[color=red]
\ssmove  0 1
\ssdropbull \ssstroke[color=red]
\ssmoveto 0 1
\ssline[color=red] 1 1 
\ssdropbull  \ssdroplabel{x} \ssstroke[color=red]
\ssmove  1 1
\ssdropbull \ssstroke[color=red]
\ssmoveto 1 4
\ssdropbull
 \ssdroplabel{z}
\ssmove  0 1
\ssdropbull 
\ssstroke[color=mygreen]
\ssmove  0 1
\ssdropbull \ssstroke[color=mygreen]
\ssmove  0 1
\ssdropbull \ssstroke[color=mygreen]
\ssmove  0 1
\ssdropbull \ssstroke[color=mygreen]
\ssmove  0 1
\ssdropbull \ssstroke[color=mygreen]
\ssmoveto 1 4 
\ssline[color=mygreen] 1 1
\ssdropbull \ssstroke[color=mygreen]
\ssmove  1 1
\ssdropbull \ssstroke[color=mygreen]
\ssmoveto 1 4
\ssline[dashed] {-1}{1} \ssarrowhead
\ssmove{1}{0}
\ssline[dashed] {-1}{1} \ssarrowhead
\ssmove{1}{0}
\ssline[dashed] {-1}{1} \ssarrowhead
\ssmove{1}{0}
\ssline[dashed] {-1}{1} \ssarrowhead
\ssmove{1}{0}
\ssline[dashed] {-1}{1} 
\ssmoveto{1}{0}
\ssdrop{\rule{0pt}{0pt}\raisebox{-60pt}{$0\!\!\pmod 8$}}
\end{sseq} \ 
\begin{sseq}[entrysize=5mm, gridstroke=.05pt,xlabelstep=1, ylabels=none]{0...2}{0...8}
\ssmove 0 1
\ssdropbull
\ssmove  0 1
\ssdropbull  \ssstroke[color=red]
\ssmove  0 1
\ssdropbull \ssstroke[color=red]
\ssmove  0 1
\ssdropbull \ssstroke[color=red]
\ssmove  0 1
\ssdropbull \ssstroke[color=red]
\ssmove  0 1
\ssdropbull \ssstroke[color=red]
\ssmove  0 1
\ssdropbull \ssstroke[color=red]
\ssmove  0 1
\ssdropbull \ssstroke[color=red]
\ssmove  0 1
\ssdropbull \ssstroke[color=red]
\ssmove  0 1
\ssdropbull \ssstroke[color=red]
\ssmoveto 0 1
\ssline[color=red] 1 1 
\ssdropbull  \ssdroplabel{x} \ssstroke[color=red]
\ssmove  1 1
\ssdropbull \ssstroke[color=red]

\ssmoveto {-1}{4}
\ssdropbull 
\ssmove  0 1
\ssdropbull \ssstroke[color=blue]
\ssmove  0 1
\ssdropbull \ssstroke[color=blue]
\ssmove  0 1
\ssdropbull \ssstroke[color=blue]
\ssmove  0 1
\ssdropbull \ssstroke[color=blue]
\ssmove  0 1
\ssdropbull \ssstroke[color=blue]

\ssmoveto {1}{5}
\ssdropbull \ssdroplabel{y}
\ssmove  0 1
\ssdropbull \ssstroke[color=blue]
\ssmove  0 1
\ssdropbull \ssstroke[color=blue]
\ssmove  0 1
\ssdropbull \ssstroke[color=blue]
\ssmove  0 1
\ssdropbull \ssstroke[color=blue]

\ssmoveto {3}{6}
\ssdropbull
\ssmove  0 1
\ssdropbull \ssstroke[color=blue]
\ssmove  0 1
\ssdropbull \ssstroke[color=blue]
\ssmove  0 1
\ssdropbull \ssstroke[color=blue]


\ssmoveto 1 5
\ssline[dashed] {-1}{1} \ssarrowhead
\ssmove{1}{0}
\ssline[dashed] {-1}{1} \ssarrowhead
\ssmove{1}{0}
\ssline[dashed] {-1}{1} \ssarrowhead
\ssmove{1}{0}
\ssline[dashed] {-1}{1}
\ssmoveto{1}{0}
\ssdrop{\rule{0pt}{0pt}\raisebox{-60pt}{$1\!\!\pmod 8$}}
\end{sseq} \ 
\begin{sseq}[entrysize=5mm, gridstroke=.05pt,xlabelstep=1, ylabels=none]{0...2}{0...8}
\ssmove  0 1
\ssdropbull
\ssmove  0 1 
\ssdropbull \ssstroke[color=red]
\ssmove  0 1
\ssdropbull \ssstroke[color=red]
\ssmove  0 1
\ssdropbull \ssstroke[color=red]
\ssmove  0 1
\ssdropbull \ssstroke[color=red]
\ssmove  0 1
\ssdropbull \ssstroke[color=red]
\ssmove  0 1
\ssdropbull \ssstroke[color=red]
\ssmove  0 1
\ssdropbull \ssstroke[color=red]
\ssmove  0 1
\ssdropbull \ssstroke[color=red]
\ssmove  0 1
\ssdropbull \ssstroke[color=red]
\ssmoveto 0 1
\ssline[color=red] 1 1 
\ssdropbull  \ssdroplabel{x} \ssstroke[color=red]
\ssmove  1 1
\ssdropbull \ssstroke[color=red]

\ssmoveto 1 6
\ssdropbull \ssdroplabel{y}
\ssmove  0 1
\ssdropbull \ssstroke[color=blue]
\ssmove  0 1
\ssdropbull \ssstroke[color=blue]
\ssmove  0 1
\ssdropbull \ssstroke[color=blue]
\ssmove  0 1
\ssdropbull \ssstroke[color=blue]

\ssmoveto {-1} {5}
\ssdropbull
\ssmove  {-1} {-1}
\ssdropbull \ssstroke[color=blue]

\ssmoveto 1 6
\ssline[dashed] {-1}{1} \ssarrowhead
\ssmove{1}{0}
\ssline[dashed] {-1}{1} \ssarrowhead
\ssmove{1}{0}
\ssline[dashed] {-1}{1} 
\ssmoveto{1}{0}
\ssdrop{\rule{0pt}{0pt}\raisebox{-60pt}{$2\!\!\pmod 8$}}

\end{sseq} \ 
\begin{sseq}[entrysize=5mm, gridstroke=.05pt, xlabelstep=1,ylabels=none]{0...2}{0...8}
\ssmove  0 1
\ssdropbull
\ssmove  0 1
\ssdropbull \ssstroke[color=red]
\ssmove  0 1
\ssdropbull \ssstroke[color=red]
\ssmove  0 1
\ssdropbull \ssstroke[color=red]
\ssmove  0 1
\ssdropbull \ssstroke[color=red]
\ssmove  0 1
\ssdropbull \ssstroke[color=red]
\ssmove  0 1
\ssdropbull \ssstroke[color=red]
\ssmove  0 1
\ssdropbull \ssstroke[color=red]
\ssmove  0 1
\ssdropbull \ssstroke[color=red]
\ssmove  0 1
\ssdropbull \ssstroke[color=red]
\ssmoveto 0 1
\ssline[color=red] 1 1 
\ssdropbull \ssdroplabel{x} \ssstroke[color=red]
\ssmove  1 1
\ssdropbull \ssstroke[color=red]

\ssmoveto 1 6
\ssdropbull
\ssdroplabel{y}
\ssmove  0 1
\ssdropbull \ssstroke[color=blue]
\ssmove  0 1
\ssdropbull \ssstroke[color=blue]
\ssmove  0 1
\ssdropbull \ssstroke[color=blue]
\ssmove  0 1
\ssdropbull \ssstroke[color=blue]


\ssmoveto {-1} {5}
\ssdropbull
\ssmove  {-1} {-1}
\ssdropbull \ssstroke[color=blue]

\ssmoveto {-1} {6}
\ssdropbull
\ssmove  {-1} {-1}
\ssdropbull \ssstroke[color=mygreen]

\ssmoveto 1 6
\ssline[dashed] {-1}{1} \ssarrowhead
\ssmove{1}{0}
\ssline[dashed] {-1}{1} \ssarrowhead
\ssmove{1}{0}
\ssline[dashed] {-1}{1} 

\ssmoveto 1 7
\ssdropbull
\ssname{green1}
\ssdroplabel{z}
\ssmove  0 1
\ssdropbull \ssname{green2} \ssstroke[color=mygreen]
\ssmove  0 1
\ssdropbull \ssstroke[color=mygreen]
\ssmove  0 1
\ssdropbull \ssstroke[color=mygreen]


\ssgoto{green1}
\ssline[dashed, curve =-.2] {-1}{1} \ssarrowhead
\ssgoto{green2}
\ssline[dashed, curve =-.2] {-1}{1} \ssarrowhead
\ssmoveto{1}{0}
\ssdrop{\rule{0pt}{0pt}\raisebox{-60pt}{$3\!\!\pmod 8$}}

\end{sseq}\
\begin{sseq}[entrysize=5mm, gridstroke=.05pt,xlabelstep=1,ylabels=none]{0...2}{0...8}
\ssmove 0 1
\ssdropbull
\ssmove  0 1
\ssdropbull  \ssstroke[color=red]
\ssmove  0 1
\ssdropbull \ssstroke[color=red]
\ssmove  0 1
\ssdropbull \ssstroke[color=red]
\ssmove  0 1
\ssdropbull \ssstroke[color=red]
\ssmove  0 1
\ssdropbull \ssstroke[color=red]
\ssmove  0 1
\ssdropbull \ssstroke[color=red]
\ssmove  0 1
\ssdropbull \ssstroke[color=red]
\ssmove  0 1
\ssdropbull \ssstroke[color=red]
\ssmoveto 0 1
\ssline[color=red] 1 1 
\ssdropbull  \ssdroplabel{x} \ssstroke[color=red]
\ssmove  1 1
\ssdropbull \ssstroke[color=red]
\ssmoveto 1 3
\ssdropbull
 \ssdroplabel{z}
\ssmove  0 1
\ssdropbull 
\ssstroke[color=mygreen]
\ssmove  0 1
\ssdropbull \ssstroke[color=mygreen]
\ssmove  0 1
\ssdropbull \ssstroke[color=mygreen]
\ssmove  0 1
\ssdropbull \ssstroke[color=mygreen]
\ssmove  0 1
\ssdropbull \ssstroke[color=mygreen]
\ssmove  0 1
\ssdropbull \ssstroke[color=mygreen]
\ssmoveto {-1}{2} 
\ssdropbull
\ssline[color=mygreen] {-1}{-1}
\ssdropbull \ssstroke[color=mygreen]

\ssmoveto 1 3
\ssline[dashed] {-1}{1} \ssarrowhead
\ssmove{1}{0}
\ssline[dashed] {-1}{1} \ssarrowhead
\ssmove{1}{0}
\ssline[dashed] {-1}{1} \ssarrowhead
\ssmove{1}{0}
\ssline[dashed] {-1}{1} \ssarrowhead
\ssmove{1}{0}
\ssline[dashed] {-1}{1} 
\ssmove{1}{0}
\ssline[dashed] {-1}{1} \ssarrowhead
\ssmove{1}{0}
\ssmoveto{1}{0}
\ssdrop{\rule{0pt}{0pt}\raisebox{-60pt}{$4\!\!\pmod 8$}}
\end{sseq} \ 
\begin{sseq}[entrysize=5mm, gridstroke=.05pt, xlabelstep=1,ylabels=none]{0...2}{0...8}
\ssmove 0 1
\ssdropbull
\ssmove  0 1
\ssdropbull  \ssstroke[color=red]
\ssmove  0 1
\ssdropbull \ssstroke[color=red]
\ssmove  0 1
\ssdropbull \ssstroke[color=red]
\ssmove  0 1
\ssdropbull \ssstroke[color=red]
\ssmove  0 1
\ssdropbull \ssstroke[color=red]
\ssmove  0 1
\ssdropbull \ssstroke[color=red]
\ssmove  0 1
\ssdropbull \ssstroke[color=red]
\ssmove  0 1
\ssdropbull \ssstroke[color=red]
\ssmove  0 1
\ssdropbull \ssstroke[color=red]
\ssmoveto 0 1
\ssline[color=red] 1 1 
\ssdropbull  \ssdroplabel{x} \ssstroke[color=red]
\ssmove  1 1
\ssdropbull \ssstroke[color=red]

\ssmoveto {-1}{2}
\ssdropbull 
\ssmove  0 1
\ssdropbull \ssstroke[color=blue]
\ssmove  0 1
\ssdropbull \ssstroke[color=blue]
\ssmove  0 1
\ssdropbull \ssstroke[color=blue]
\ssmove  0 1
\ssdropbull \ssstroke[color=blue]
\ssmove  0 1
\ssdropbull \ssstroke[color=blue]
\ssmove  0 1
\ssdropbull \ssstroke[color=blue]
\ssmove  0 1
\ssdropbull \ssstroke[color=blue]

\ssmoveto {1}{3}
\ssdropbull \ssdroplabel{y}
\ssmove  0 1
\ssdropbull \ssstroke[color=blue]
\ssmove  0 1
\ssdropbull \ssstroke[color=blue]
\ssmove  0 1
\ssdropbull \ssstroke[color=blue]
\ssmove  0 1
\ssdropbull \ssstroke[color=blue]
\ssmove  0 1
\ssdropbull \ssstroke[color=blue]
\ssmove  0 1
\ssdropbull \ssstroke[color=blue]

\ssmoveto {3}{6}
\ssdropbull
\ssmove  0 1
\ssdropbull \ssstroke[color=blue]
\ssmove  0 1
\ssdropbull \ssstroke[color=blue]
\ssmove  0 1
\ssdropbull \ssstroke[color=blue]


\ssmoveto 1 3
\ssline[dashed] {-1}{1} \ssarrowhead
\ssmove{1}{0}
\ssline[dashed] {-1}{1} \ssarrowhead
\ssmove{1}{0}
\ssline[dashed] {-1}{1} \ssarrowhead
\ssmove{1}{0}
\ssline[dashed] {-1}{1} \ssarrowhead
\ssmove{1}{0}
\ssline[dashed] {-1}{1} \ssarrowhead
\ssmove{1}{0}
\ssline[dashed] {-1}{1}
\ssmoveto{1}{0}
\ssdrop{\rule{0pt}{0pt}\raisebox{-60pt}{$5\!\!\pmod 8$}}
\end{sseq} \ 
\begin{sseq}[entrysize=5mm, gridstroke=.05pt, xlabelstep=1,ylabels=none]{0...2}{0...8}
\ssmove  0 1
\ssdropbull
\ssmove  0 1 
\ssdropbull \ssstroke[color=red]
\ssmove  0 1
\ssdropbull \ssstroke[color=red]
\ssmove  0 1
\ssdropbull \ssstroke[color=red]
\ssmove  0 1
\ssdropbull \ssstroke[color=red]
\ssmove  0 1
\ssdropbull \ssstroke[color=red]
\ssmove  0 1
\ssdropbull \ssstroke[color=red]
\ssmove  0 1
\ssdropbull \ssstroke[color=red]
\ssmove  0 1
\ssdropbull \ssstroke[color=red]
\ssmove  0 1
\ssdropbull \ssstroke[color=red]
\ssmoveto 0 1
\ssline[color=red] 1 1 
\ssdropbull  \ssdroplabel{x} \ssstroke[color=red]
\ssmove  1 1
\ssdropbull \ssstroke[color=red]

\ssmoveto 1 3
\ssdropbull \ssdroplabel{y}
\ssmove  0 1
\ssdropbull \ssstroke[color=blue]
\ssmove  0 1
\ssdropbull \ssstroke[color=blue]
\ssmove  0 1
\ssdropbull \ssstroke[color=blue]
\ssmove  0 1
\ssdropbull \ssstroke[color=blue]
\ssmove  0 1
\ssdropbull \ssstroke[color=blue]
\ssmove  0 1
\ssdropbull \ssstroke[color=blue]
\ssmove  0 1
\ssdropbull \ssstroke[color=blue]

\ssmoveto {1} {3}
\ssline[color=blue, dotted]  {1} {1}
\ssdropbull
\ssmove  1 1
\ssdropbull \ssstroke[color=blue]

\ssmoveto 1 3
\ssline[dashed] {-1}{1} \ssarrowhead
\ssmove{1}{0}
\ssline[dashed] {-1}{1} \ssarrowhead
\ssmove{1}{0}
\ssline[dashed] {-1}{1} \ssarrowhead
\ssmove{1}{0}
\ssline[dashed] {-1}{1} \ssarrowhead
\ssmove{1}{0}
\ssline[dashed] {-1}{1} \ssarrowhead
\ssmove{1}{0}
\ssline[dashed] {-1}{1} 
\ssmoveto{1}{0}
\ssdrop{\rule{0pt}{0pt}\raisebox{-60pt}{$6\!\!\pmod 8$}}

\end{sseq} \ 
\begin{sseq}[entrysize=5mm, gridstroke=.05pt,xlabelstep=1,ylabels=none]{0...2}{0...8}
\ssmove  0 1
\ssdropbull
\ssdroplabel{y}
\ssmove  0 1
\ssdropbull \ssstroke[color=red]
\ssmove  0 1
\ssdropbull \ssstroke[color=red]
\ssmove  0 1
\ssdropbull \ssstroke[color=red]
\ssmove  0 1
\ssdropbull \ssstroke[color=red]
\ssmove  0 1
\ssdropbull \ssstroke[color=red]
\ssmove  0 1
\ssdropbull \ssstroke[color=red]
\ssmove  0 1
\ssdropbull \ssstroke[color=red]
\ssmove  0 1
\ssdropbull \ssstroke[color=red]
\ssmove  0 1
\ssdropbull \ssstroke[color=red]
\ssmoveto 0 1
\ssline[color=red] 1 1 
\ssdropbull 
\ssdroplabel{x}
\ssstroke[color=red]
\ssmove  1 1
\ssdropbull \ssstroke[color=red]


\ssmoveto  1 3
\ssdropbull
\ssname{13blue}
\ssdroplabel{y}
\ssmove  0 1
\ssdropbull \ssname{14blue}\ssstroke[color=blue]
\ssmove  0 1
\ssdropbull \ssname{15blue}\ssstroke[color=blue]
\ssmove  0 1
\ssdropbull \ssname{16blue}\ssstroke[color=blue]
\ssmove  0 1
\ssdropbull \ssname{17blue} \ssstroke[color=blue]
\ssmove  0 1
\ssdropbull \ssname{18blue} \ssstroke[color=blue]
\ssmove  0 1
\ssdropbull \ssstroke[color=blue]

\ssmoveto 1 3 
\ssline[color=blue,dotted] 1 1 
\ssdropbull 
\ssmove  1 1
\ssdropbull \ssstroke[color=blue]


\ssmoveto  1 4
\ssdropbull
\ssname{14green}
\ssdroplabel{z}
\ssmove  0 1
\ssdropbull \ssname{15green} \ssstroke[color=mygreen]
\ssmove  0 1
\ssdropbull \ssname{16green} \ssstroke[color=mygreen]
\ssmove  0 1
\ssdropbull \ssname{17green} \ssstroke[color=mygreen]
\ssmove  0 1
\ssdropbull \ssname{18green} \ssstroke[color=mygreen]
\ssmove  0 1
\ssdropbull \ssstroke[color=mygreen]

\ssgoto {13blue}
\ssline[dashed] {-1}{1} \ssarrowhead
\ssgoto {14blue}
\ssline[dashed] {-1}{1} \ssarrowhead
\ssgoto {15blue}
\ssline[dashed] {-1}{1} \ssarrowhead
\ssgoto {16blue}
\ssline[dashed] {-1}{1} 
\ssgoto {17blue}
\ssline[dashed] {-1}{1} 
\ssgoto {18blue}
\ssline[dashed] {-1}{1} 

\ssgoto {14green}
\ssline[color=mygreen] 1 1 
\ssdropbull \ssstroke[color=mygreen]
\ssmove  1 1
\ssdropbull \ssstroke[color=mygreen]

\ssgoto{14green}
\ssline[dashed, curve =-.2] {-1}{1} \ssarrowhead
\ssgoto{15green}
\ssline[dashed, curve =-.2] {-1}{1} \ssarrowhead
\ssgoto{16green}
\ssline[dashed, curve =-.2] {-1}{1} \ssarrowhead
\ssgoto{17green}
\ssline[dashed, curve =-.2] {-1}{1} \ssarrowhead
\ssgoto{18green}
\ssline[dashed, curve =-.2] {-1}{1} \ssarrowhead

\ssmoveto{1}{0}
\ssdrop{\rule{0pt}{0pt}\raisebox{-60pt}{$7\!\!\pmod 8$}}

\end{sseq}\\
% $0\!\!\pmod 8$\hspace{0.4cm}
% $1\!\!\pmod 8$\hspace{0.4cm}
% $2\!\!\pmod 8$\hspace{0.4cm}
% $3\!\!\pmod 8$\hspace{0.4cm}
% $4\!\!\pmod 8$\hspace{0.4cm}
% $5\!\!\pmod 8$\hspace{0.4cm}
% $6\!\!\pmod 8$\hspace{0.4cm}
% $7\!\!\pmod 8$
\caption{The $\Gilt$-filtered $E_2$ pages near the column $i=1$.}
\label{fig:n07}
\end{figure}


The $\Gilt$ filtration produces a spectral sequence going from this
associated graded to the $E_2$ page of the Adams spectral sequence for
$X=\RP^{-n}$. Let us prove that the differentials are given by the
dashed arrows in Figure~\ref{fig:n07}. For $n$ even, we can appeal to
Equation~\eqref{eq:even-support}, which says that the $E_2$ page of
the Adams spectral sequence (that is, what is left after canceling
generators with the differentials) should be supported in the
half-plane $s \leq \frac{i}{2} + \frac{n}{2}+1.$ This means that the
differentials from either the blue or the green pattern to the red
pattern exist.

For $n$ odd, Equation~\eqref{eq:odd-support} applies only to the even numbered columns such as $i=0$. This
shows that the top part of the red pattern must be canceled by the
differentials, but this does not suffice to determine the
differentials. An additional argument is provided by considering the
inclusion maps $\iota: \RP^{n} \to \RP^{n+1}$
with duals
\[
  D\iota\from \RP^{-n-1} \to \RP^{-n}.
\]
These induce maps between the spectral sequences (and their associated
graded). The maps induced by $D\iota$ always take the red pattern
isomorphically to the red pattern, because this comes from the common
piece $G_0$. Moreover, when the consecutive values of $n$ both have a
green pattern in the column $i=1$, these are taken to each other
injectively; the same goes for the blue patterns. (Even when one blue
pattern comes from $R_1$ and the other from $R_2$, these can be
related by filtering $R_2$ to break it into a copy of $R_1$ and one of
$\FF_2[3]$, as before.)  Since the map $D\iota$ commutes with the
differential, and since we know the differentials for $n$ even, we
deduce that for $n$ odd they are also as drawn in
Figure~\ref{fig:n07}.


The map we are interested in is 
\[
  (D\pi)_*\from ko_1(\RP^{-n}) \otimes \ZZ_2 \to ko_1(S^{-n}) \otimes \ZZ_2
\]
for $n\equiv 0,1,3,7\!\!\pmod 8$, cf.~Equation~\eqref{eq:ns}.  The
cofiber sequence
$S^n \xrightarrow{\pi} \RP^{n} \xrightarrow{\iota} \RP^{n+1} $ has
dual
\[
  \RP^{-n-1} \xrightarrow{\iota}  \RP^{-n} \xrightarrow{D\pi} S^{-n}.
\]
producing a long exact sequence on $ko$-homology. Hence, the induced
maps on $ko_1$ satisfy $(D\pi)_* \circ (D\iota)_* = 0.$ Thus, an
element in the image of $(D\iota)_*$ is automatically taken to $0$ by
$(D\pi)_*$. This is the case with the elements labeled $x$ in
Figure~\ref{fig:n07}. These survive to the $E_2$ page of the Adams
spectral sequence, and they are the image under $(D\iota)_*$ of the
corresponding $x$ for $n+1$.

In the cases $n \equiv 0,1\!\!\pmod 8$, we see that $x$ is the only surviving element in the column $i=1$. Thus, we have $ ko_1(\RP^{-n})  \otimes \ZZ_2 \cong \ZZ/2$, and the map 
\[
  (D\pi)_*\from \ZZ/2 \to \ZZ/2
\]
is zero.

In the cases $n \equiv 3,7\!\!\pmod 8$, we again have that $x$ maps to $0$ under $(D\pi)_*$. However, we also have to consider the elements $h_0^r(z + h_0 y)$ for $r \geq 0$, which survive after the action of the differentials. They are mapped by $(D\pi)_*$ to something in $ko_1(S^{-n})$, i.e., in the green patterns in Figure~\ref{fig:n07}. 

Observe that $\pi\from S^n \to \RP^n$ induces the zero map on the reduced homology and cohomology with $\FF_2$ coefficients; hence, so does $D\pi$. This shows that the map induced by $D\pi$ on the $E_2$ page of the Adams spectral sequence,
\[
  (D\pi)_*\from \Ext^{*,*}_{\Aone}(H^*(\RP^{-n};\FF_2),\FF_2) \to \Ext^{*,*}_{\Aone}(H^*(S^{-n};\FF_2),\FF_2)
\]
is trivial. This implies that $(D\pi)_*$ also vanishes on the $E_{\infty}$ page, but it may not vanish on $ko_1$, because of extension problems. Nevertheless, we deduce that $(D\pi)_*$ increases the Adams filtration degree $s$, and therefore takes $z + h_0 y$ to an element of higher $s$ in the green pattern. All such elements are in the image of $h_0$. Since $h_0$ corresponds to multiplication by $2$ in $ko_1(S^{-n})\otimes \ZZ_2 \cong \ZZ_2$, we conclude that the image of 
$$  (D\pi)_*\from ko_1(\RP^{-n}) \otimes \ZZ_2 \to ko_1(S^{-n}) \otimes \ZZ_2$$
is $2$-divisible. 

This completes the proof when $n \geq 4$. It remains to analyze the
cases $n=1$ and $n=3$. When $n=1$, we have $\RP^{-1} \cong
S^{-1}$. Then, the map $D\pi\from S^{-1} \to S^{-1}$ is of degree $2$,
and induces the zero map on $ko_{1}(S^{-1}) \cong \ZZ/2$.


\begin{figure}
   \centering
\begin{sseq}[entrysize=5mm, gridstroke=.05pt]{-3...3}{0...6}
\ssmoveto {-3}{0}
\ssdropbull
\ssmove  0 1
\ssdropbull \ssstroke[color=mygreen]
\ssmove  0 1
\ssdropbull \ssstroke[color=mygreen]
\ssmove  0 1
\ssdropbull \ssstroke[color=mygreen]
\ssmove  0 1
\ssdropbull \ssstroke[color=mygreen]
\ssmove  0 1
\ssdropbull \ssstroke[color=mygreen]
\ssmove  0 1
\ssdropbull \ssstroke[color=mygreen]
\ssmove  0 1
\ssdropbull \ssstroke[color=mygreen]
\ssmove  0 1
\ssdropbull \ssstroke[color=mygreen]
\ssmove  0 1
\ssdropbull \ssstroke[color=mygreen]
\ssmove{0}{1}
\ssdropbull \ssstroke[color=mygreen]
\ssmove{0}{-10}
\ssline[color=mygreen]{1}{1}
\ssdropbull
\ssmove 1 1 
\ssdropbull \ssstroke[color=mygreen]
\ssmove 2 1
\ssdropbull \ssdroplabel{z}
\ssmove  0 1
\ssdropbull \ssstroke[color=mygreen]
\ssmove  0 1
\ssdropbull \ssstroke[color=mygreen]
\ssmove  0 1
\ssdropbull \ssstroke[color=mygreen]
\ssmove  0 1
\ssdropbull \ssstroke[color=mygreen]
\ssmove  0 1
\ssdropbull \ssstroke[color=mygreen]
\ssmove{0}{1}
\ssdropbull \ssstroke[color=mygreen]

\ssmoveto{-2}{0}
\ssdropbull
\ssmove  1 1
\ssdropbull \ssstroke[color=red]
\ssmove  1 1
\ssdropbull \ssstroke[color=red]
\ssmove  {0} {-1}
\ssdropbull \ssstroke[color=red]
\ssmove  1 1
\ssdropbull \ssdroplabel{x} \ssstroke[color=red]
\ssmove  1 1
\ssdropbull \ssstroke[color=red]

\end{sseq}
\caption{The associated graded of the $E_2$ page for $n=3$.}
\label{fig:RP3}
\end{figure}

When $n=3$, either the $\Filt$ or the $\Gilt$ filtration splits
$H^*(\RP^{-3}; \FF_2)$ into a copy of $\Azero[-2]$ and one of
$\FF_2[-3]$. Thus, the $E_2$ page of the Adams spectral sequence has
an associated graded as shown in Figure~\ref{fig:RP3}.  The proof that
the image of $(D\pi)_*$ is $2$-divisible is similar to the case
$n=8j+3$ for $j \geq 1$. The element $x$ is the image of the
corresponding $x$ for $\RP^{-4}$, and therefore maps to $0$ under
$(D\pi)_*$. The element $z$ maps to $0$ on the $E_{\infty}$ page of
the Adams spectral sequence. Therefore, on $ko_1$ it maps with
something in the green pattern with a higher value of $s$, so its
image must be $2$-divisible.
\end{proof}

\bibliographystyle{amsplain}
\bibliography{ko}



\end{document}

%%% Local Variables:
%%% mode: latex
%%% TeX-master: t
%%% End:
